problem:  
Let $\pi: E_k \to S^4$ denote the smooth $S^3$–bundle associated to the principal $SU(2)$–bundle over $S^4$ with second Chern class $c_2=k\in\mathbb{Z}$.  Each connection $A$ on the principal bundle defines a **connection metric** $g_{A,r}$ on $E_k$ by giving the fibers the round metric of radius $r$ and the horizontal subspaces determined by $A$.  

Known examples (Bérard‑Bergery, Grove–Ziller) show that for certain small integers $k$, some $E_k$ admit positively curved metrics, sometimes homogeneous or biquotient.  However, no classification is known at the level of **connection metrics**.

**Question:**  
For the canonical base $(S^4,h_\text{round})$, determine for which integers $k$ there exists a smooth connection $A$ and scaling factor $r>0$ such that the connection metric $g_{A,r}$ on $E_k$ has strictly positive sectional curvature.  
Equivalently, does the space $\mathcal C_+(S^4,\!SU(2),k)=\{(A,r): K_{g_{A,r}}>0\}$ depend nontrivially on $k$, possibly vanishing for large $|k|$?  

Why is it a "good" problem:  
This problem aims to sharply classify the interplay between **topological twisting (Chern number)** and **Riemannian positivity** for one of the most fundamental bundle geometries: $S^3$–bundles over $S^4$.  Unlike broad existence questions, it isolates a concrete topological parameter and asks for its geometric threshold, potentially revealing whether positive sectional curvature can occur only within a finite range of bundle classes—or whether curvature positivity is topologically unconstrained in this setting.  

It connects concrete computations (O’Neill’s formulas, curvature of $SU(2)$–connections) with deep topological invariants (the integer $k$), bridging global Riemannian geometry, algebraic topology, and gauge theory.  Any classification or obstruction discovered here would illuminate the precise geometric content of the integer class $k$ and might explain why positive curvature has so far appeared only for specific bundles.  

The statement is extremely simple—merely asking “which $k$ allow a positively curved connection metric on $E_k \to S^4$?”—yet its resolution would require new insight into the quantitative relation between bundle topology and curvature.  Its simplicity, cross‑disciplinary nature, and proximity to one of the central mechanisms of bundle geometry make it a paradigmatic “good” research‑level question.